\documentclass[twoside,10pt]{article}

\usepackage[T1]{fontenc}
\usepackage{lmodern}
\usepackage{amsmath,amsthm,amssymb,amsfonts}
\usepackage{mathtools}
\usepackage{array}
\usepackage[small,labelfont=bf,labelsep=period]{caption}
\usepackage[
  letterpaper,
  total={5.5in,8in},
  vmarginratio=1:1,
  hmarginratio=1:1,
  headsep=21pt
]{geometry}
\usepackage{xcolor}
\usepackage[
  breaklinks=false,
  linktocpage=true
]{hyperref}
\hypersetup{
  colorlinks=true,
  linkcolor={purple},
  urlcolor={purple},
  citecolor={blue!80!black},
  pdfauthor={Leslie P. Polzer},
  pdftitle={Matchgate Commutator Graphs as Cover Graphs of Rectangular Minuscule Lattices: Brownian-Bridge Distances, Graph Complexity, and a Krylov Bound},
  pdfsubject={Brownian-bridge distances and complexity in matchgate commutator-graph components},
  pdfkeywords={matchgate commutator graph, rectangular minuscule lattice, token graph, Wiener index, Brownian bridge, Wasserstein distance, graph complexity, Krylov complexity}
}

\medmuskip=4mu
\thickmuskip=8mu
\setlength{\emergencystretch}{2em}

\numberwithin{equation}{section}

\newtheoremstyle{slbody}
  {8pt}
  {8pt}
  {\normalfont\slshape}
  {}
  {\bfseries}
  {.}
  {0.5em}
  {\thmname{#1}\thmnumber{ #2}%
   \thmnote{ {\normalfont\itshape(\ignorespaces#3\unskip)}}}
\theoremstyle{slbody}
\newtheorem{theorem}{Theorem}[section]
\newtheorem{proposition}[theorem]{Proposition}
\newtheorem{lemma}[theorem]{Lemma}
\newtheorem{corollary}[theorem]{Corollary}
\theoremstyle{definition}
\newtheorem{definition}[theorem]{Definition}
\theoremstyle{remark}
\newtheorem{remark}[theorem]{Remark}

\renewcommand{\qedsymbol}{$\blacksquare$}

\newcommand{\E}{\mathbb{E}}
\newcommand{\Pp}{\mathbb{P}}
\newcommand{\Var}{\operatorname{Var}}
\newcommand{\1}{\mathbf{1}}
\newcommand{\N}{\mathbb{N}}

\title{Matchgate commutator graphs as cover graphs of rectangular minuscule lattices:\\
Brownian-bridge distances, graph complexity, and a Krylov bound}
\author{Leslie P. Polzer}
\newcommand{\authoraffiliation}{Independent Researcher}
\newcommand{\authoremail}{polzer@fastmail.com}
\date{September 4, 2026}

\makeatletter
\renewcommand{\section}{%
  \@startsection{section}{1}{\z@}%
    {-3.5ex \@plus -1ex \@minus -.2ex}%
    {2.3ex \@plus .2ex}%
    {\normalfont\large\bfseries}}
\renewcommand{\@seccntformat}[1]{%
  \csname the#1\endcsname.\enspace}

\def\tagform@#1{%
  \maketag@@@{(\ignorespaces{\oldstylenums{#1}}\unskip\@@italiccorr)}}
\renewcommand{\eqref}[1]{%
  \textup{{\normalfont(\oldstylenums{\ref{#1}}}\normalfont)}}

\newcommand{\shorttitle}{MATCHGATE GRAPHS AND MINUSCULE LATTICES}
\newcommand{\shortauthors}{LESLIE P. POLZER}
\newcommand{\ps@bookheader}{%
  \renewcommand\@oddfoot{\hfil}%
  \renewcommand\@evenfoot{\hfil}%
  \renewcommand\@oddhead{%
    \ifnum\value{page}>1
      {\small\hfil\shortauthors}\hfil\thepage
    \else
      \hfil
    \fi}%
  \renewcommand\@evenhead{%
    \thepage\hfil\small\shorttitle\hfil}}
\pagestyle{bookheader}

\renewcommand{\maketitle}{%
  \begin{center}
    {\large\bfseries\@title}\\
    \vskip36pt
    \ifx\@author\@empty
      \mbox{}\par
    \else
      {\scshape\@author}\par
      \medskip
      {\small\authoraffiliation}\par
      \smallskip
      {\small\texttt{\authoremail}}\par
      \smallskip
      {\small\@date}\par
    \fi
    \vskip36pt
  \end{center}}
\makeatother

\renewenvironment{abstract}
  {\quotation\noindent\small{\bfseries\abstractname.\enspace}}
  {\endquotation}

\begin{document}
\raggedbottom
\maketitle

\begin{abstract}
Matchgate commutator-graph components are isomorphic, as unweighted graphs, to
the cover graphs of rectangular minuscule lattices.  Defant, F\'eray, Nadeau,
and Williams previously proved
the fixed-aspect-ratio Brownian-bridge distance law for these lattices, with
moment convergence, and obtained a closed Wiener-index
formula~\cite{DefantEtAl}.  Combining that result with the
matchgate--minuscule-lattice identification shows that Conjecture~14 of West
et al.~\cite[Conjecture~14]{WestEtAl} holds.  We make this implication
explicit and give an independent finite-population derivation valid for every
sequence $k_n/(2n)\to\rho\in(0,1)$.  We also derive a graph-complexity asymptotic after
averaging jointly over a uniform initial vertex and an independent Haar
matchgate unitary, and prove a separate pointwise graph--Krylov comparison.
If $S$ and $T$ are independent uniform $k_n$-subsets of
$[2n]$, write
$D_{2n,k_n}(S,T)=\sum_i|s_i-t_i|$.  Then
\[
 \frac{D_{2n,k_n}}{n^{3/2}}
 \Rightarrow 4\sqrt{\rho(1-\rho)}\int_0^1|B(t)|\,dt,
 \qquad
 \E D_{2n,k_n}\sim
 \frac{\sqrt{2\pi}}2\sqrt{\rho(1-\rho)}\,n^{3/2},
\]
where $B$ is the standard Brownian bridge with covariance
$\min\{s,t\}-st$.  In the central (half-filled) component $k_n=n$, the
coefficient multiplying the Brownian area is $2$ and the mean constant is
$\sqrt{2\pi}/4$.  We also record the variance asymptotic, extract the leading
finite-size correction from the published closed form, and give a finite sum
obtained by conditioning on the symmetric-difference size.  The Krylov
comparison is for generator-supported time-independent Hamiltonians; no
Haar-averaged Krylov asymptotic is asserted.

\vskip5pt
\noindent\textbf{Keywords.}\enspace Matchgate commutator graph; rectangular
minuscule lattice; token graph; Wiener index; Brownian bridge; Wasserstein
distance; graph complexity; Krylov complexity.

\vskip5pt
\noindent\textbf{2020 Mathematics Subject Classification.}\enspace
60F17; 05C12; 06A07; 60J65; 81P68.
\end{abstract}

\vskip50pt
\baselineskip=13.5pt

\section{Introduction}

The commutator graph was introduced by Diaz et
al.~\cite[Supplemental Definition~6]{DiazEtAl}.
West et al.~\cite[Section~IV and Appendix~C]{WestEtAl} describe the component
$\mathcal C_\kappa$ of the matchgate commutator graph by the $\kappa$-element
subsets of $[2n]$.  Their
Appendix~C, Eq.~(C2), identifies the graph distance between
$S=(s_1<\cdots<s_\kappa)$ and $T=(t_1<\cdots<t_\kappa)$ as
\begin{equation}\label{eq:matching-distance}
 d_{\mathcal C_\kappa}(S,T)=\sum_{i=1}^{\kappa}|s_i-t_i|.
\end{equation}
Theorem~13 of West et al. gives a finite sum for the average of
\eqref{eq:matching-distance}, and Conjecture~14 asserts that the central value
is superlinear and subquadratic in $n$.
We retain $\kappa$ for the matchgate component label when discussing West et
al.; elsewhere $k$ denotes the same subset cardinality, so $\kappa=k$ when
$N=2n$.

The same graph and distance were studied earlier in the language of minuscule
lattices.  Let
\[
 \mathcal P_{m,\ell}:=J([m]\times[\ell])
\]
be the distributive lattice of order ideals of an $m\times\ell$ rectangle.
Reading the membership word of a $k$-subset of $[N]$ as an up--down path
identifies the token graph with the undirected cover graph (Hasse graph) of
$\mathcal P_{k,N-k}$.  Put $C_S(r):=|S\cap[r]|$ and define $C_T(r)$
similarly.  If $p_r$ and $q_r$ are the heights of the paths for $S$ and $T$
after $r$ steps, then $p_r=2C_S(r)-r$ and $q_r=2C_T(r)-r$.  Consequently,
the path-area distance in Eq.~(2) of Defant et al. is
\[
 \frac12\sum_{r=1}^{N-1}|p_r-q_r|
 =\sum_{r=1}^{N-1}|C_S(r)-C_T(r)|
 =D_{N,k}(S,T),
\]
where the last equality is also proved directly in
Lemma~\ref{lem:transport}.  Defant, F\'eray, Nadeau, and
Williams~\cite[Corollary~3 and Proposition~5]{DefantEtAl} give a
closed formula for its Wiener index and prove the fixed-aspect-ratio
Brownian-bridge limit in distribution and in moments.  Defant et al. write
$d(G)$ for the distance sum over all ordered vertex pairs.  Under the standard
unordered-pair convention for the Wiener index $W(G)$, symmetry and the zero
diagonal give $d(G)=2W(G)$; hence the average used by West et al. is
$d(G)/|V(G)|^2=2W(G)/|V(G)|^2$.

Defant et al.~\cite[Proposition~5]{DefantEtAl} state the result for rectangles
$\mathcal P_{\lfloor\alpha\ell\rfloor,\ell}$ along a fixed ray, with
$\alpha>0$ fixed.  Along this ray, the subset density tends to
$\rho:=\alpha/(1+\alpha)$.  Writing the total path length as
$N=k+\ell$ gives $\ell\sim(1-\rho)N$; on the matchgate subsequence $N=2n$,
this is $\ell\sim2(1-\rho)n$.  Thus their
$\ell^{3/2}$-normalized limit has, in the present $n^{3/2}$ normalization,
coefficient
\[
 [2(1-\rho)]^{3/2}\sqrt{2\alpha(1+\alpha)}
 =4\sqrt{\rho(1-\rho)}.
\]
Their moment convergence gives the same mean and variance constants; at
$\alpha=1$ it gives the central case.  Thus their theorem, combined with the
matchgate component identification of West et al., already implies
Conjecture~14.  The present finite-population proof gives a triangular-array
formulation in which the two rectangle dimensions may vary arbitrarily subject
only to
$k_n/(2n)\to\rho$, equivalently
$k_n/(2n-k_n)\to\rho/(1-\rho)$, and supplies direct moment bounds.  It also
derives the Haar matchgate specialization, gives a
projection proof of the graph--Krylov comparison, and obtains a central finite
sum by conditioning on the symmetric-difference size.  The Haar statement is
a joint expectation over an independent uniform initial Pauli vertex and Haar
matchgate unitary; it is not a pointwise-in-vertex scaling statement.

Table~\ref{tab:provenance} separates the prior results used here from the
deductions and alternative arguments given in this paper.

\begin{table}[hbp]
\centering
\small
\begin{tabular}{>{\raggedright\arraybackslash}p{0.34\linewidth}
                >{\raggedright\arraybackslash}p{0.58\linewidth}}
Result & Source and role in this paper\\
\hline
Central Brownian-area law and moments
& Proposition~5 of Defant et al.~\cite[Proposition~5]{DefantEtAl}, combined
with the matchgate-component identification of West et
al.~\cite[Section~IV and Appendix~C]{WestEtAl}; recorded here as a matchgate
consequence.\\
Arbitrary density sequences
& Finite-population proof here for every $k_n/(2n)\to\rho$, giving a
density-sequence formulation and direct second- and fourth-moment bounds.\\
Exact finite mean
& Corollary~3 of Defant et al.~\cite[Corollary~3]{DefantEtAl}; equivalently
Theorem~4.3 of Bourn and Erickson~\cite[Theorem~4.3]{BournErickson}.\\
Central finite-size formulas
& The $n^{1/2}$ correction is extracted here from the published mean; the
conditional bridge-and-gap sum is derived here.\\
Graph and Krylov complexity
& Matchgate specialization of Proposition~7 of West et
al.~\cite[Proposition~7 and Lemma~9]{WestEtAl}; the statement of their
Lemma~9 is proved here by a tail-projection argument that avoids its
cross-term step.\\
\end{tabular}
\caption{Provenance of the principal results and the role of the present
arguments.}
\label{tab:provenance}
\end{table}

The alternative probabilistic proof uses one-dimensional transport.  The
matching distance is the area between two cumulative counting functions, and
each centered counting function is a finite-population bridge.  Ros\'en's
finite-population invariance principle~\cite[Theorem~13.1]{Rosen} supplies
convergence in $C[0,1]$.  Proposition~\ref{prop:second-moment} proves the
uniform integrability needed for convergence of the first two moments directly
from hypergeometric moment bounds.

The token-graph terminology is due to Fabila-Monroy et
al.~\cite{FabilaMonroy}; the same graph is also called the $k$th symmetric power of a
path~\cite{AudenaertEtAl}.  The ordered-coordinate distance formula appears,
for example, in Ibarra and Rivera~\cite[Lemma~5.2]{IbarraRivera}.  The transport
identity below is the discrete one-dimensional $1$-Wasserstein identity.  In
one dimension, $W_1$ is the area between cumulative distribution functions
\cite{Vallender,DeAngelisGray}.  Vallender's later addendum corrects Remark~1
and formula~(3) of the original note; it does not alter the CDF-area identity
used here~\cite{VallenderAddendum}.

Bourn and Willenbring~\cite{BournWillenbring} study expected one-dimensional
earth mover distance for uniform weak compositions and its probability-simplex
limit.  The stars-and-bars encoding identifies a uniform $k$-subset of $[N]$
with a uniform
weak composition of $k$ into $N-k+1$ parts, and Lemma~\ref{lem:transport}
identifies the corresponding earth mover distance with $D_{N,k}$.  Bourn and
Erickson~\cite[Theorem~4.3]{BournErickson} use the Type~A minuscule-lattice
Wiener formula of Defant et al. to obtain the discrete expected-distance
formula; the substitutions $s=k$ and $b=N-k+1$ give exactly
Proposition~\ref{prop:closed-form}.  They also interpret the distance through
symmetric differences of Young diagrams.  Proposition~\ref{prop:central-exact}
below instead conditions on $|S\setminus T|$ and decomposes the distance into a
simple-walk bridge area and exchangeable gaps.

Brownian-bridge limits for $L^1$ empirical-process and Wasserstein functionals
are classical; see del
Barrio, Gin\'e, and Matr\'an~\cite{delBarrioGineMatran} and their correction to
separate heavy-tail statements~\cite{delBarrioCorrection}.  Those corrected
statements are not used here.

Janson~\cite[Section~20]{Janson} surveys the Brownian-bridge area and its
moments, while Schwerdtfeger~\cite[Section~12.2]{Schwerdtfeger} proves joint
moment convergence for the absolute area and endpoint of a Bernoulli walk.
Section~6 records the published rectangular-lattice closed form, gives an
independent bridge-and-gap identity for its central specialization, and
extracts the leading $n^{1/2}$ finite-size correction.  Section~5 supplies the
explicit Haar specialization and a projection proof of the
Hamiltonian-dependent graph--Krylov comparison.

\section{The matchgate metric and discrete transport}

For $N\geq1$, write $[N]=\{1,\ldots,N\}$ and $[0]=\varnothing$, and let
$P_N$ be the path graph with vertex set $[N]$ and edges
$\{r,r+1\}$ for $1\leq r<N$.  For integers $1\leq k<N$, let
$\Omega_{N,k}$ be the set of $k$-subsets of $[N]$.  Write each
$S\in\Omega_{N,k}$ as
$S=(s_1<\cdots<s_k)$ and define
\[
 D_{N,k}(S,T):=\sum_{i=1}^k|s_i-t_i|,
 \qquad
 C_S(r):=|S\cap[r]|.
\]

When $N=2n$, define the Jordan--Wigner Majorana operators
\[
 c_{2j-1}=Z_1\cdots Z_{j-1}X_j,\qquad
 c_{2j}=Z_1\cdots Z_{j-1}Y_j\qquad(1\leq j\leq n),
\]
and put
\begin{equation}\label{eq:majorana-support-convention}
 g_a:=-ic_ac_{a+1}\quad(1\leq a<2n),
 \qquad
 P_S:=i^{k(k-1)/2}c_{s_1}\cdots c_{s_k}
 \quad(S\in\Omega_{2n,k}).
\end{equation}

\begin{proposition}[Configuration-graph distance]\label{prop:graph-metric}
Consider the graph on $\Omega_{N,k}$ in which one particle may move by one
lattice step to an unoccupied neighboring site.  This is the $k$-token graph
$F_k(P_N)$ and the Hasse graph of the rectangular minuscule lattice
$\mathcal P_{k,N-k}$, and its graph distance is $D_{N,k}$.  When $N=2n$, the
map $S\mapsto P_S$ in \eqref{eq:majorana-support-convention} identifies this
graph with the degree-$k$ matchgate commutator-graph component
$\mathcal C_k$ generated by $\{g_a:1\leq a<2n\}$.  In particular, its metric
is \eqref{eq:matching-distance}.
\end{proposition}

\begin{proof}
For fixed $T$, every edge changes $D_{N,k}(\,\cdot\,,T)$ by at most one, so
every path from $S$ to $T$ has length at least $D_{N,k}(S,T)$.  For the reverse
inequality, suppose first
that $s_i<t_i$ for some $i$, and choose the largest such $i$.  The move
$s_i\mapsto s_i+1$ stays in $[N]$: if $i<k$, then
$s_i+1\leq s_{i+1}\leq N$, while if $i=k$, then
$s_k+1\leq t_k\leq N$.  It is also unblocked: if $i<k$ and
$s_i+1=s_{i+1}$, maximality gives $s_{i+1}\geq t_{i+1}$, whereas
$s_{i+1}=s_i+1\leq t_i<t_{i+1}$, a contradiction.  Thus this move reduces
$D_{N,k}$ by one.  If no such $i$ exists but $S\neq T$, choose the smallest
$i$ for which $s_i>t_i$.  The move stays in $[N]$: if $i>1$, then
$s_i-1\geq s_{i-1}\geq1$, while if $i=1$, then
$s_1-1\geq t_1\geq1$.  If $i>1$ and the move were blocked, then
$s_i-1=s_{i-1}$; minimality would give $s_{i-1}\leq t_{i-1}$, whereas
$s_{i-1}=s_i-1\geq t_i>t_{i-1}$, a contradiction.  Hence
$s_i\mapsto s_i-1$ is an allowed move and again reduces $D_{N,k}$ by one.
Iteration constructs a path of length $D_{N,k}$.

For the lattice identification, encode $S$ by a path with an up-step in
position $r$ when $r\in S$ and a down-step otherwise.  Such paths bound the
order ideals of a $k\times(N-k)$ rectangle.  A token move swaps consecutive
up--down steps, which adds or removes one cell, so token edges are precisely
Hasse edges.  This is the path model of Defant et
al.~\cite[Section~1.3 and Eq.~(2)]{DefantEtAl}.  The token-graph terminology
is introduced by Fabila-Monroy et al.~\cite{FabilaMonroy}; the definition and
the same distance formula are also recorded by Ibarra and
Rivera~\cite[Definition in Section~1 and Lemma~5.2]{IbarraRivera}.  For
$N=2n$, the Majoranas satisfy $c_a^\dagger=c_a$ and
$\{c_a,c_b\}=2\delta_{ab}I$.  Direct multiplication gives
\[
 g_{2j-1}=Z_j\quad(1\leq j\leq n),\qquad
 g_{2j}=X_jX_{j+1}\quad(1\leq j<n),
\]
so the $g_a$ are precisely the matchgate generator family.  Reversing the
ordered Majorana product introduces the sign
$(-1)^{k(k-1)/2}$.  If $e=k(k-1)/2$, then
\[
 P_S^\dagger=(-i)^e(-1)^e c_{s_1}\cdots c_{s_k}=P_S,
\]
so the chosen phase makes every $P_S$ a Hermitian Pauli string.  If
$A_a=\{a,a+1\}$, the canonical anticommutation
relations give the sign below: commuting the two Majorana factors indexed by
$A_a$ through the $k$ factors indexed by $S$ requires
$2k-|A_a\cap S|$ anticommutations.
\[
 g_aP_S=(-1)^{|A_a\cap S|}P_Sg_a.
\]
Thus $[g_a,P_S]\neq0$ precisely when $S$ contains exactly one of $a,a+1$.
In that case $g_aP_S$ is a nonzero scalar multiple of
$P_{S\mathbin{\triangle}A_a}$.  On supports this replaces the occupied member
of the adjacent pair by the unoccupied member, which is exactly one token
move.  For example, when $n=2$ and $k=2$, the support $\{1,3\}$ is joined
through $g_1$ to $\{2,3\}$ and through $g_3$ to $\{1,4\}$.  The graph is
connected by the distance-decreasing construction above, so these Pauli
strings form the component $\mathcal C_k$.  This recovers Eqs.~(50)--(51) and
the adjacency rule of West et
al.~\cite[Section~IV and Appendix~C]{WestEtAl}.
\end{proof}

For the empirical measures
\[
 \mu_S=\frac1k\sum_{i=1}^k\delta_{s_i},
 \qquad
 \mu_T=\frac1k\sum_{i=1}^k\delta_{t_i},
\]
one has
$D_{N,k}(S,T)=kW_1(\mu_S,\mu_T)$, where $W_1$ uses the ground cost
$|x-y|$ on $[N]$.  Lemma~\ref{lem:transport} is therefore the
lattice form of the one-dimensional identity expressing $W_1$ as the area
between the two cumulative distribution functions
\cite{Vallender,DeAngelisGray}.

\begin{lemma}[Discrete transport identity]\label{lem:transport}
For all $S,T\in\Omega_{N,k}$,
\begin{equation}\label{eq:transport}
 D_{N,k}(S,T)=\sum_{r=1}^{N-1}|C_S(r)-C_T(r)|.
\end{equation}
\end{lemma}

\begin{proof}
For $a,b\in[N]$,
\[
 |a-b|=\sum_{r=1}^{N-1}|\1_{\{a\leq r\}}-\1_{\{b\leq r\}}|.
\]
Fix $r$.  The nonzero differences
$\1_{\{s_i\leq r\}}-\1_{\{t_i\leq r\}}$ cannot have both signs.  If a
positive difference occurred at $i$ and a negative difference at $j$, then
$i<j$ would imply $t_i<t_j\leq r<t_i$, while $j<i$ would imply
$s_j<s_i\leq r<s_j$; both are impossible.  Hence the sum of the absolute
values at level $r$ equals
$|C_S(r)-C_T(r)|$.  Summing first over $i$ and then over $r$ proves
\eqref{eq:transport}.
\end{proof}

\section{Finite-population bridges and moment control}

Let $S_N,T_N$ be independent and uniform in $\Omega_{N,k_N}$, put
$p_N=k_N/N$, and let $G_N^S$ be the polygonal process whose values on the
grid are
\begin{equation}\label{eq:single-bridge}
 G_N^S(m/N)=\frac{C_{S_N}(m)-mp_N}{\sqrt N},
 \qquad 0\leq m\leq N.
\end{equation}

\begin{lemma}[Finite-population bridge limit]\label{lem:bridge-fclt}
If $p_N\to\rho\in(0,1)$, then in $C[0,1]$ with the uniform norm,
\[
 G_N^S\Rightarrow\sqrt{\rho(1-\rho)}B,
\]
where $B$ is the standard Brownian bridge with
$\operatorname{Cov}(B(s),B(t))=\min\{s,t\}-st$.  If $G_N^T$ is the
independent copy formed from $T_N$, then
\begin{equation}\label{eq:difference-bridge}
 H_N:=G_N^S-G_N^T
 \Rightarrow \sqrt{2\rho(1-\rho)}B.
\end{equation}
\end{lemma}

\begin{proof}
Apply Ros\'en~\cite[Theorem~13.1, pp.~408--411]{Rosen} to the population
containing $k_N$ ones and $N-k_N$ zeros, exposed in uniformly random order
through all $N$ positions.  Its mean is $p_N$, its population
variance in Ros\'en's normalization is
\[
 \sigma_N^2=\frac{N}{N-1}p_N(1-p_N)
 \longrightarrow\rho(1-\rho).
\]
For the full sample of size $N$, Ros\'en's normalized polygon has, at time
$m/N$, the value
\[
 \frac{C_{S_N}(m)-mp_N}{\sigma_N\sqrt N}
 =\frac{G_N^S(m/N)}{\sigma_N}.
\]
Ros\'en's Noether condition is satisfied because, with
$Q_N^2=\sum_{i=1}^N(a_i-p_N)^2=Np_N(1-p_N)$,
\[
 \frac{\max_i|a_i-p_N|}{Q_N}
 \leq\frac{1}{\sqrt{Np_N(1-p_N)}}\longrightarrow0.
\]
The sample size and population size are both $N$, so the sampling fraction in
Ros\'en's theorem is $\lambda=1$ and his limiting process $W(1,1)$ is the
standard Brownian bridge.  The cited theorem therefore applies to exactly this
polygon and gives $G_N^S/\sigma_N\Rightarrow B$ in $C[0,1]$.  Since
$\sigma_N\to\sqrt{\rho(1-\rho)}$, Slutsky's theorem proves the first assertion.
Independence makes the joint law of $(G_N^S,G_N^T)$ the product of its
marginal laws, so it converges to two independent bridges.  Subtraction is
continuous on $C[0,1]^2$, and the difference of two independent standard
bridges has the law of $\sqrt2B$.  This proves
\eqref{eq:difference-bridge}.
\end{proof}

\begin{proposition}[Second- and fourth-moment bounds]\label{prop:second-moment}
For independent uniform $S,T\in\Omega_{N,k}$,
\begin{equation}\label{eq:second-moment-bound}
 \E D_{N,k}(S,T)^2
 \leq \frac{k(N-k)(N^2-1)}{3N}.
\end{equation}
Moreover,
\begin{equation}\label{eq:fourth-moment-bound}
 \E D_{N,k}(S,T)^4\leq N^6.
\end{equation}
In particular, for every sequence $1\leq k_n<2n$,
\[
 \Var(D_{2n,k_n})=O(n^3).
\]
Both families
\[
 \left\{\frac{D_{2n,k_n}}{n^{3/2}}:n\geq1\right\},
 \qquad
 \left\{\left(\frac{D_{2n,k_n}}{n^{3/2}}\right)^2:n\geq1\right\}
\]
are uniformly integrable.
\end{proposition}

\begin{proof}
Let $\Delta_r=C_S(r)-C_T(r)$ and put $p=k/N$.  For the membership
indicators $I_j=\1_{\{j\in S\}}$, one has
$\E I_j=p$, $\Var(I_j)=p(1-p)$, and, for $i\neq j$,
\[
 \operatorname{Cov}(I_i,I_j)
 =\frac{k(k-1)}{N(N-1)}-p^2
 =-\frac{p(1-p)}{N-1}.
\]
Summing these covariances over the first $r$ sites, and then using the
independence of $S$ and $T$, gives
\begin{equation}\label{eq:hypergeom-var}
 \E\Delta_r^2
 =2\Var(C_S(r))
 =\frac{2k(N-k)}{N-1}\frac rN\left(1-\frac rN\right).
\end{equation}
By Lemma~\ref{lem:transport} and Cauchy--Schwarz,
\[
 \E D_{N,k}^2
 \leq (N-1)\sum_{r=1}^{N-1}\E\Delta_r^2
 =2k(N-k)\sum_{r=1}^{N-1}\frac{r(N-r)}{N^2}.
\]
The finite sum is $(N^2-1)/(6N)$, which proves
\eqref{eq:second-moment-bound}.  The variance bound follows from
$\Var(D_{N,k})\leq\E D_{N,k}^2$.  If $N=2n$ and $k=k_n$, then
$k_n(2n-k_n)\leq n^2$, so division by $n^3$ gives
\[
 \sup_n \E\left(\frac{D_{2n,k_n}}{n^{3/2}}\right)^2<\infty.
\]
An $L^2$-bounded family is uniformly integrable.

For the fourth moment, put $q_r=\min(r,N-r)$.  If $X$ is Bernoulli with
mean $p$, define
\[
 \psi(\lambda):=\log\E\exp(\lambda(X-p)).
\]
Then
$\psi(0)=\psi'(0)=0$ and
$\psi''(\lambda)=p_\lambda(1-p_\lambda)\leq1/4$, where $p_\lambda$ is the
success probability under exponential tilting.  Taylor's formula therefore
gives $\psi(\lambda)\leq\lambda^2/8$ for every $\lambda\in\mathbb R$.
Hoeffding's convex comparison for sampling without replacement
\cite[Theorem~4]{Hoeffding}, in the form stated by
Ros\'en~\cite[Theorem~3.1]{Rosen}, applied separately to the two independent
counts, gives for $r\leq N/2$
\[
 \E e^{\lambda(C_S(r)-rp)}\leq e^{r\lambda^2/8},
 \qquad
 \E e^{-\lambda(C_T(r)-rp)}\leq e^{r\lambda^2/8}.
\]
Independence therefore gives, with $q_r=\min(r,N-r)$,
\[
 \E e^{\lambda\Delta_r}\leq e^{q_r\lambda^2/4}.
\]
For $r>N/2$ this uses
$\Delta_r=-[(k-C_S(r))-(k-C_T(r))]$ and samples the complementary block of
size $N-r$.  For $x>0$, minimizing
$-\lambda x+q_r\lambda^2/4$ at $\lambda=2x/q_r$, and then applying the same
bound to $-\Delta_r$, gives
\[
 \Pp(|\Delta_r|\geq x)\leq2e^{-x^2/q_r}.
\]
Integration of this tail gives
\[
 \E|\Delta_r|^4
 =4\int_0^\infty x^3\Pp(|\Delta_r|>x)\,dx
 \leq8\int_0^\infty x^3e^{-x^2/q_r}\,dx
 =4q_r^2.
\]
H\"older's inequality and $\sum_{r=1}^{N-1}q_r^2\leq N^3/4$ now yield
\[
 \E D_{N,k}^4
 \leq(N-1)^3\sum_{r=1}^{N-1}\E|\Delta_r|^4
 \leq N^6.
\]
In fact, \eqref{eq:fourth-moment-bound} gives
\[
 \sup_{n\geq1}\sup_{1\leq k<2n}
 \E\left(\frac{D_{2n,k}}{n^{3/2}}\right)^4\leq64.
\]
Thus the squared normalized distances are bounded in $L^2$, which proves the
last uniform integrability assertion.
\end{proof}

\section{The Brownian-area law}

Set
\[
 \mathcal A:=\int_0^1|B(t)|\,dt.
\]

\begin{theorem}[Density-sequence matchgate distance]\label{thm:main}
Let $N=2n$, let $k_n/(2n)\to\rho\in(0,1)$, and let $S_n,T_n$ be independent
uniform $k_n$-subsets of $[2n]$.  Then
\begin{equation}\label{eq:distribution-main}
 \frac{D_{2n,k_n}(S_n,T_n)}{n^{3/2}}
 \Rightarrow 4\sqrt{\rho(1-\rho)}\,\mathcal A,
\end{equation}
and
\begin{equation}\label{eq:expectation-main}
 \E D_{2n,k_n}(S_n,T_n)
 =\frac{\sqrt{2\pi}}2\sqrt{\rho(1-\rho)}\,n^{3/2}
   +o(n^{3/2}).
\end{equation}
Its variance satisfies
\begin{equation}\label{eq:variance-main}
 \Var(D_{2n,k_n})
 =\rho(1-\rho)\left(\frac{28}{15}-\frac\pi2\right)n^3
  +o(n^3).
\end{equation}
For the central component $k_n=n$ this specializes to
\begin{equation}\label{eq:central-main}
 \frac{D_{2n,n}}{n^{3/2}}\Rightarrow2\mathcal A,
 \qquad
 \E D_{2n,n}=\frac{\sqrt{2\pi}}4n^{3/2}+o(n^{3/2}).
\end{equation}
In that case
\begin{equation}\label{eq:central-variance}
 \Var(D_{2n,n})
 =\left(\frac7{15}-\frac\pi8\right)n^3+o(n^3).
\end{equation}
\end{theorem}

\begin{proof}
With $N=2n$, form $H_N$ from $S_n,T_n$ as in
Lemma~\ref{lem:bridge-fclt}.  At the grid points it satisfies
$H_N(r/N)=(C_{S_n}(r)-C_{T_n}(r))/\sqrt N$.  Thus
Lemma~\ref{lem:transport} gives
\begin{equation}\label{eq:riemann-area}
 \frac{D_{N,k_n}}{N^{3/2}}
 =\frac1N\sum_{r=1}^{N-1}|H_N(r/N)|.
\end{equation}
The missing right-endpoint term is zero because
$C_{S_n}(N)=C_{T_n}(N)=k_n$.  On each grid interval, $H_N$ is linear and its
endpoint values differ by at most $N^{-1/2}$.  Since the absolute-value map is
one-Lipschitz,
\begin{equation}\label{eq:riemann-error}
 \left|
 \frac1N\sum_{r=1}^{N-1}|H_N(r/N)|
 -\int_0^1|H_N(t)|\,dt
 \right|\leq N^{-1/2}.
\end{equation}
The map $f\mapsto\int_0^1|f(t)|\,dt$ is one-Lipschitz on $C[0,1]$ in the
uniform norm.  Lemma~\ref{lem:bridge-fclt}, the continuous mapping theorem,
and \eqref{eq:riemann-error} therefore yield
\[
 \frac{D_{N,k_n}}{N^{3/2}}
 \Rightarrow\sqrt{2\rho(1-\rho)}\,\mathcal A.
\]
Since $N^{3/2}=2\sqrt2\,n^{3/2}$, this is
\eqref{eq:distribution-main}.

Proposition~\ref{prop:second-moment} makes
$D_{2n,k_n}/n^{3/2}$ uniformly integrable, so its expectations converge to
the expectation of the limit in \eqref{eq:distribution-main}.  For every
$t\in[0,1]$, $B(t)$ is centered Gaussian with variance $t(1-t)$; hence
Tonelli's theorem gives
\begin{align*}
 \E\mathcal A
 &=\int_0^1\E|B(t)|\,dt
 =\sqrt{\frac2\pi}\int_0^1\sqrt{t(1-t)}\,dt \\
 &=\sqrt{\frac2\pi}\,\mathrm B(3/2,3/2)
 =\sqrt{\frac2\pi}\,\frac\pi8
 =\frac{\sqrt{2\pi}}8.
\end{align*}
Multiplication by $4\sqrt{\rho(1-\rho)}$ proves
\eqref{eq:expectation-main}; setting $\rho=1/2$ proves
\eqref{eq:central-main}.

Put $X_n=D_{2n,k_n}/n^{3/2}$ and
$X=4\sqrt{\rho(1-\rho)}\,\mathcal A$.  The fourth-moment bound in
Proposition~\ref{prop:second-moment} makes $\{X_n^2\}$ uniformly integrable,
while \eqref{eq:distribution-main} gives $X_n\Rightarrow X$.  The continuous
mapping theorem therefore gives $X_n^2\Rightarrow X^2$, and uniform
integrability yields $\E X_n^2\to\E X^2$.  Together with the already proved
convergence $\E X_n\to\E X$, this gives $\Var(X_n)\to\Var(X)$.  The
Brownian-bridge area has
\begin{equation}\label{eq:brownian-area-second-moment}
 \E\mathcal A^2=\frac7{60};
\end{equation}
this is recorded, with the same normalization, in
Janson~\cite[Section~20, Eq.~(135) and Table~2, p.~107]{Janson}.  The variance of the limit in
\eqref{eq:distribution-main} is therefore
\[
 16\rho(1-\rho)\frac7{60}
 -\left(\frac{\sqrt{2\pi}}2\sqrt{\rho(1-\rho)}\right)^2
 =\rho(1-\rho)\left(\frac{28}{15}-\frac\pi2\right).
\]
This proves \eqref{eq:variance-main}; setting $\rho=1/2$ gives
\eqref{eq:central-variance}.
\end{proof}

\begin{corollary}[Conjecture 14 as a corollary]\label{cor:conjecture}
Let $A(n,n)$ be the central-component average in Theorem~13 of
West et al.~\cite[Theorem~13 and Conjecture~14]{WestEtAl}.  Then
\[
 A(n,n)=\frac{\sqrt{2\pi}}4n^{3/2}+o(n^{3/2}).
\]
Consequently $A(n,n)=\omega(n)$ and $A(n,n)=o(n^2)$, as asserted in their
Conjecture~14.
\end{corollary}

\begin{proof}
The average $A(n,n)$ is $\E D_{2n,n}$ by Proposition~\ref{prop:graph-metric}
and the definition preceding Theorem~13 of West et al.  The conclusion is
the central case of Theorem~\ref{thm:main}.  It also follows from
Defant et al.~\cite[Proposition~5]{DefantEtAl} at rectangle aspect ratio one;
the present corollary records that published lattice result in the matchgate
notation of West et al.
\end{proof}

\section{Graph and Krylov complexity}

Let
\begin{equation}\label{eq:matchgate-generators}
 \Gamma_n:=\{g_a:1\leq a<2n\}
 =\{Z_j:1\leq j\leq n\}
 \cup\{X_jX_{j+1}:1\leq j<n\},
\end{equation}
and let $\mathcal M_n$ be the closed subgroup of the unitary group generated
by $e^{-i\theta g}$ with $g\in\Gamma_n$ and $\theta\in\mathbb R$.  As a
closed subgroup of the finite-dimensional unitary group, $\mathcal M_n$ is
compact.  It is the matchgate group used by West et
al.~\cite[Table~I and the discussion following Eq.~(44)]{WestEtAl}; its
dynamical Lie algebra is isomorphic to $\mathfrak{so}(2n)$, and its commutator
graph has the components $\mathcal C_k$ identified in
Proposition~\ref{prop:graph-metric}.

Equip the operator space with the normalized Hilbert--Schmidt inner product
\begin{equation}\label{eq:normalized-HS}
 \langle A,B\rangle_{\mathrm{HS}}:=2^{-n}\operatorname{Tr}(A^\dagger B),
\end{equation}
so that the Pauli strings are orthonormal.  If
$O=\sum_{Q\in\mathcal C_k}c_QQ$, define graph complexity based at
$P\in\mathcal C_k$ by
\begin{equation}\label{eq:graph-complexity-definition}
 \mathsf G_P(O):=\sum_{Q\in\mathcal C_k}
 d_{\mathcal C_k}(P,Q)|c_Q|^2.
\end{equation}
This is the normalization of West et al.~\cite[Eq.~(37)]{WestEtAl}.

\begin{lemma}[Uniform Haar coefficient]\label{lem:haar-coefficient}
For every $1\leq k<2n$, every $P,Q\in\mathcal C_k$, and Haar-distributed
$U\in\mathcal M_n$,
\[
 \E_U\left|\langle Q,U^\dagger P U\rangle_{\mathrm{HS}}\right|^2
 =\frac1{|\mathcal C_k|}
 =\binom{2n}{k}^{-1}.
\]
\end{lemma}

\begin{proof}
Use the Majorana operators, generators, and canonical Pauli representatives
from \eqref{eq:majorana-support-convention}.  For $1\leq a<2n$,
$V_a:=e^{-i\pi g_a/4}$ satisfies
\[
 V_a^\dagger c_aV_a=-c_{a+1},\qquad
 V_a^\dagger c_{a+1}V_a=c_a,\qquad
 V_a^\dagger c_bV_a=c_b\quad(b\notin\{a,a+1\}).
\]
More generally, the canonical anticommutation relations give
\begin{align*}
 e^{i\theta g_a}c_ae^{-i\theta g_a}
 &=\cos(2\theta)c_a-\sin(2\theta)c_{a+1},\\
 e^{i\theta g_a}c_{a+1}e^{-i\theta g_a}
 &=\sin(2\theta)c_a+\cos(2\theta)c_{a+1}.
\end{align*}
Thus conjugation by $\mathcal M_n$ preserves each degree-$k$ Majorana span.
The induced actions of the $V_a$ on supports are the adjacent transpositions,
which generate the full symmetric group on the $2n$ Majorana labels.  Given
two $k$-element supports, choose a permutation carrying
the first support to the second and express it as a product of adjacent
transpositions.  The corresponding product of the $V_a$ maps the canonical
Pauli representative of the first support to that of the second up to a sign.
Thus the action is transitive on the one-dimensional basis subspaces
$\mathbb R P_S$, equivalently on the support labels, although conjugation may
introduce a sign.  This transitivity is the only representation-theoretic
input below; the argument does not assume that the degree-$k$ Majorana
representation is irreducible.

For fixed $P$ and $Q\in\mathcal C_k$, put
\[
 a_Q:=\E_U\left|\langle Q,U^\dagger P U\rangle_{\mathrm{HS}}\right|^2.
\]
Given $Q,R\in\mathcal C_k$, choose $V\in\mathcal M_n$ and
$\varepsilon\in\{-1,1\}$ such that $V^\dagger QV=\varepsilon R$.  Unitary
invariance of the Hilbert--Schmidt
inner product and right invariance of Haar measure, with $W=UV^\dagger$,
give
\begin{align*}
 a_R
 &=\E_U\left|\langle V^\dagger QV,U^\dagger P U\rangle_{\mathrm{HS}}\right|^2
   \qquad (\varepsilon=\pm1)\\
 &=\E_U\left|\langle Q,VU^\dagger PUV^\dagger\rangle_{\mathrm{HS}}\right|^2
 =\E_W\left|\langle Q,W^\dagger P W\rangle_{\mathrm{HS}}\right|^2
 =a_Q.
\end{align*}
Hence $a_Q$ is independent of $Q$.  Conjugation by $\mathcal M_n$ preserves
the degree-$k$ Majorana span, and Parseval's identity gives, for every $U$,
\[
 \sum_{Q\in\mathcal C_k}
 \left|\langle Q,U^\dagger P U\rangle_{\mathrm{HS}}\right|^2
 =\|U^\dagger P U\|_{\mathrm{HS}}^2=1.
\]
Taking expectations proves $|\mathcal C_k|a_Q=1$.  The component size is
$|\mathcal C_k|=\binom{2n}{k}$ by the Majorana-subset identification.
This is the matchgate specialization of Proposition~7 of West et
al.~\cite[Proposition~7 and Eq.~(36)]{WestEtAl}, with the normalized pairing
made explicit.
\end{proof}

\begin{corollary}[Matchgate graph complexity under vertex and Haar averaging]
\label{cor:complexity}
Let $k_n/(2n)\to\rho\in(0,1)$, let $P_n$ be uniform in the matchgate component
$\mathcal C_{k_n}$, and let $U_n$ be independently Haar distributed in the
group $\mathcal M_n$.  Then
\[
 \E_{P_n,U_n}\mathsf G_{P_n}(U_n^\dagger P_nU_n)
 =\frac{\sqrt{2\pi}}2\sqrt{\rho(1-\rho)}\,n^{3/2}
  +o(n^{3/2}).
\]
\end{corollary}

\begin{proof}
Fix $P_n$ and, for $Q\in\mathcal C_{k_n}$, write the normalized Pauli
coefficient as
\[
 c_Q(U):=\langle Q,U^\dagger P_nU\rangle_{\mathrm{HS}}
 =2^{-n}\operatorname{Tr}(Q^\dagger U^\dagger P_nU).
\]
Lemma~\ref{lem:haar-coefficient} gives
$\E_U|c_Q(U)|^2=|\mathcal C_{k_n}|^{-1}$.  Substitution in
\eqref{eq:graph-complexity-definition} gives Eq.~(38) of West et
al.~\cite[Eq.~(38)]{WestEtAl}:
\[
 \E_U\mathsf G_{P_n}(U^\dagger P_nU)
 =\frac{1}{|\mathcal C_{k_n}|}
   \sum_{Q\in\mathcal C_{k_n}}d_{\mathcal C_{k_n}}(P_n,Q).
\]
This fixed-vertex mean need not be constant in $P_n$.  Averaging also over
$P_n$ gives exactly $\E D_{2n,k_n}$, so
Theorem~\ref{thm:main} gives the stated asymptotic.
\end{proof}

We use the operator-Krylov/Lanczos convention of Parker et
al.~\cite[Section~III, Eqs.~(3)--(8), and Section~V.A,
Eq.~(26)]{ParkerEtAl}.

\begin{corollary}[Graph complexity is bounded by Krylov complexity]\label{cor:krylov}
Fix $n$, an initial Pauli string $P\in\mathcal C_k$, and a time-independent
Hamiltonian
\begin{equation}\label{eq:matchgate-Hamiltonian}
 H=\sum_{g\in\Gamma_n}h_g g,\qquad h_g\in\mathbb R.
\end{equation}
Thus $H$ lies in the real linear span of the nearest-neighbor generator family
$\Gamma_n$; this restriction is part of the statement.
Put
\[
 P_t=e^{iHt}Pe^{-iHt},
\]
and apply the Lanczos construction to the Liouvillian
$\mathcal L_H=[H,\,\cdot\,]$ using \eqref{eq:normalized-HS}, terminating at
the first zero residual.  Write
\[
 d_K:=\dim\operatorname{span}\{\mathcal L_H^\ell P:\ell\geq0\}
\]
and denote the resulting orthonormal basis by
$O_0=P,O_1,\ldots,O_{d_K-1}$.  If
\[
 P_t=\sum_{j=0}^{d_K-1}d_j(t)O_j,
 \qquad
 \mathsf K_{H,P}(t):=\sum_{j=0}^{d_K-1}j|d_j(t)|^2,
\]
then, for every $t\in\mathbb R$,
\[
 \mathsf G_P(P_t)\leq \mathsf K_{H,P}(t).
\]
Consequently, for every probability law on triples $(P,H,t)$ supported on the
stated class, $\E\mathsf G_P(P_t)\leq\E\mathsf K_{H,P}(t)$.
\end{corollary}

\begin{proof}
Since the Pauli generators are Hermitian and $h_g\in\mathbb R$, one has
$H^\dagger=H$.  The Liouvillian is self-adjoint for the normalized
Hilbert--Schmidt inner product: cyclicity of the trace gives
\begin{align*}
 \langle A,\mathcal L_HB\rangle_{\mathrm{HS}}
 &=2^{-n}\operatorname{Tr}(A^\dagger HB-A^\dagger BH)\\
 &=2^{-n}\operatorname{Tr}(A^\dagger HB-HA^\dagger B)
 =\langle\mathcal L_HA,B\rangle_{\mathrm{HS}}.
\end{align*}
For $j\geq0$, let
\[
 \mathcal F_j:=\operatorname{span}\{Q\in\mathcal C_k:
 d_{\mathcal C_k}(P,Q)\leq j\},
 \qquad
 \mathcal K_j:=\operatorname{span}
 \{O_0,\ldots,O_{\min\{j,d_K-1\}}\}.
\]
For Pauli strings $g$ and $Q$, one has $[g,Q]=0$ when they commute and
$[g,Q]=2gQ$ when they anticommute.  If $g\in\Gamma_n$, the latter product is,
up to phase, the Pauli string obtained by moving one occupied Majorana label
across the adjacent pair defining $g$, and hence is adjacent to $Q$ in the
commutator graph.  Consequently
\[
 \mathcal L_H\mathcal F_j\subseteq\mathcal F_{j+1},
 \qquad \mathcal L_H^jP\in\mathcal F_j.
\]
The Lanczos construction for the self-adjoint operator $\mathcal L_H$ gives
$O_j\in\operatorname{span}\{\mathcal L_H^\ell P:0\leq\ell\leq j\}$
before its first zero residual, hence
$\mathcal K_j\subseteq\mathcal F_j$ for every $j\geq0$.  Its completed basis
spans the cyclic space $\operatorname{span}\{\mathcal L_H^\ell P:\ell\geq0\}$,
which is invariant under $\mathcal L_H$ in the finite-dimensional operator space.
Therefore $P_t=e^{it\mathcal L_H}P$ belongs to $\mathcal K_{d_K-1}$.
Let $\Pi_j^{\mathcal F}$ and $\Pi_j^{\mathcal K}$ be the orthogonal projections
onto these two spaces.  Expanding the integer weights in
\eqref{eq:graph-complexity-definition} as sums of unit increments gives
\begin{equation}\label{eq:graph-tail-projections}
 \mathsf G_P(P_t)
 =\sum_{j\geq1}\left\|(I-\Pi_{j-1}^{\mathcal F})P_t\right\|_{\mathrm{HS}}^2.
\end{equation}
The analogous tail-sum identity for the Lanczos coefficients is
\begin{equation}\label{eq:krylov-tail-projections}
 \mathsf K_{H,P}(t)
 =\sum_{j\geq1}\left\|(I-\Pi_{j-1}^{\mathcal K})P_t\right\|_{\mathrm{HS}}^2.
\end{equation}
Both sums have only finitely many nonzero terms: for $j\geq d_K$, one has
$P_t\in\mathcal K_{j-1}\subseteq\mathcal F_{j-1}$.
Because $\mathcal K_{j-1}\subseteq\mathcal F_{j-1}$, orthogonal projection
onto $\mathcal F_{j-1}$ leaves a residual of no greater norm than projection
onto $\mathcal K_{j-1}$.  Comparing
\eqref{eq:graph-tail-projections} and
\eqref{eq:krylov-tail-projections} term by term proves the pointwise
inequality.  Thus the argument supplies a self-contained projection proof of
the statement recorded as Lemma~9 by West et
al.~\cite[Lemma~9 and Eqs.~(A78)--(A86)]{WestEtAl}.  Integrating
the pointwise inequality proves the expectation statement.
\end{proof}

\begin{remark}[Relation to Lemma~9 of West et al.]
The inference from Eq.~(A81) to Eq.~(A82) in the appendix proof of Lemma~9 in
West et al.~\cite[Eqs.~(A81)--(A82)]{WestEtAl} is not justified by the stated
support condition alone.  The summands there contain cross terms in Lanczos
coefficients and Pauli
amplitudes, which need not be nonnegative real numbers, so the bound
$\ell(p,q)\leq\min\{n,n'\}$ cannot be applied to each summand as an order
inequality.  Already for real scalars, coordinatewise inequalities
$a_i\leq b_i$ do not imply $\sum_i a_iz_i\leq\sum_i b_iz_i$ when the $z_i$
have mixed signs: take $a=(1,0)$, $b=(1,1)$, and $z=(1,-2)$.  Equations~
\eqref{eq:graph-tail-projections} and
\eqref{eq:krylov-tail-projections} instead express both complexities as sums
of squared tail norms.  Those summands are nonnegative, and
$\mathcal K_{j-1}\subseteq\mathcal F_{j-1}$ gives their term-by-term
comparison.  Under the generator-supported hypothesis
$H\in\operatorname{span}_{\mathbb R}\Gamma_n$ in
\eqref{eq:matchgate-Hamiltonian}, the preceding proof therefore establishes
the cited inequality without the cross-term step.
\end{remark}

\begin{remark}[Why the Haar and Krylov statements are separate]
The Haar-distributed endpoint $U^\dagger P U$ in
Corollary~\ref{cor:complexity} does not by itself specify a Hamiltonian, a
Liouvillian, or a Lanczos basis.  Krylov complexity is therefore not a function
of that endpoint alone.  Corollary~\ref{cor:krylov} applies once the
Hamiltonian evolution supplying those data has been specified; no
Haar-averaged Krylov quantity is asserted without such a joint dynamical law.
\end{remark}

\section{Exact finite-size evaluation}

\subsection{Exact rectangular-lattice mean}

The rectangular-minuscule-lattice Wiener index gives a closed form for every
component.  In discrete earth-mover notation, the same expression is
recorded by Bourn and Erickson~\cite[Theorem~4.3]{BournErickson}: setting their total
mass to $s=k$ and their number of bins to $b=N-k+1$ gives the formula below.

\begin{proposition}[Exact rectangular-lattice mean (Defant et al.)]\label{prop:closed-form}
For independent uniform $S,T\in\Omega_{N,k}$,
\begin{equation}\label{eq:closed-form}
 \E D_{N,k}(S,T)
 =\frac{k(N-k)}{4N+2}
   \frac{\binom{2N+2}{2k+1}}{\binom Nk^2}.
\end{equation}
In particular,
\begin{equation}\label{eq:central-closed-form}
 \E D_{2n,n}
 =\frac{n^2}{8n+2}
   \frac{\binom{4n+2}{2n+1}}{\binom{2n}{n}^2}.
\end{equation}
\end{proposition}

\begin{proof}
Proposition~\ref{prop:graph-metric} identifies the graph with the Hasse graph
of $\mathcal P_{k,N-k}$, which has $\binom Nk$ vertices.  Defant et
al.~\cite[Corollary~3]{DefantEtAl} define their Wiener index as the distance
sum over all ordered vertex pairs and compute
\[
 \sum_{S,T\in\Omega_{N,k}}D_{N,k}(S,T)
 =\frac{k(N-k)}{4N+2}\binom{2N+2}{2k+1}.
\]
Division by $\binom Nk^2$ proves \eqref{eq:closed-form}; setting $N=2n$ and
$k=n$ gives \eqref{eq:central-closed-form}.
\end{proof}

\subsection{Central bridge-and-gap identity}

The Young-diagram interpretation of Bourn and Erickson expresses the distance
as a symmetric difference.  The central component also has the following
separate conditional decomposition, obtained from the symmetric-difference
size $M=|S\setminus T|$.  Let $W^{(m)}$ be a
uniform simple random-walk bridge of length $2m$, with increments $+1$ and
$-1$ occurring $m$ times each, and put
\[
 \mathcal A_m=\sum_{j=1}^{2m-1}|W^{(m)}_j|.
\]

\begin{lemma}[Mean area of a simple-walk bridge]\label{lem:discrete-bridge-area}
For every $m\geq1$,
\begin{equation}\label{eq:discrete-bridge-area}
 \E\mathcal A_m=\frac{m4^m}{2\binom{2m}{m}}.
\end{equation}
\end{lemma}

\begin{proof}
Let $L_{m,h}$ be the number of pairs $(w,j)$ in which $w$ is a bridge of
length $2m$, $0\leq j\leq2m$, and $w_j=h$.  The generating function for
unrestricted walks from $0$ to $h$ is
\[
 P_h(z)=\sum_{j\geq0}\#\{w:0\to h\text{ in }j\text{ steps}\}z^j
 =\frac{r(z)^{|h|}}{\sqrt{1-4z^2}},
 \qquad
 r(z)=\frac{1-\sqrt{1-4z^2}}{2z}.
\]
This follows by solving the recurrence
$P_h=\1_{\{h=0\}}+z(P_{h-1}+P_{h+1})$ subject to
$P_{-h}=P_h$ and $P_h(z)\in z^{|h|}\mathbb R[[z^2]]$.  These conditions select
the displayed formal-power-series solution uniquely.  Splitting a bridge at
its marked time and reversing the segment
after that time gives
$\sum_{m\geq0}L_{m,h}z^{2m}=P_h(z)^2$.  Therefore, if $T_m$ is the total
absolute area over all $\binom{2m}{m}$ bridges, then
the marked times $j=0$ and $j=2m$ contribute zero, so their inclusion in
$L_{m,h}$ agrees with the interior-time sum defining $\mathcal A_m$.  Hence
\begin{align*}
 \sum_{m\geq0}T_mz^{2m}
 &=\sum_{h\in\mathbb Z}|h|P_h(z)^2
 =\frac{2}{1-4z^2}\sum_{h\geq1}h r(z)^{2h}\\
 &=\frac{2r(z)^2}{(1-4z^2)(1-r(z)^2)^2}
 =\frac{2z^2}{(1-4z^2)^2}.
\end{align*}
In the last equality we used $r=z(1+r^2)$.  The coefficient of $z^{2m}$
is $m4^m/2$.  Division by the number of bridges proves
\eqref{eq:discrete-bridge-area}.
\end{proof}

\begin{proposition}[Central symmetric-difference identity]\label{prop:central-exact}
For independent uniform $n$-subsets $S,T$ of $[2n]$,
\begin{equation}\label{eq:central-exact}
 \E D_{2n,n}
 =\frac{2n+1}{2\binom{2n}{n}}
 \sum_{m=1}^n
 \binom{n}{m}^2
 \frac{m4^m}{(2m+1)\binom{2m}{m}}.
\end{equation}
\end{proposition}

\begin{proof}
Let $M=|S\setminus T|=|T\setminus S|$.  Conditional on $S$, choosing $T$
gives
\begin{equation}\label{eq:M-law}
 \Pp(M=m)=\frac{\binom{n}{m}^2}{\binom{2n}{n}},
 \qquad 0\leq m\leq n.
\end{equation}
Conditional on $M=m$, mark the $2m$ sites in the symmetric difference by
$+1$ on $S\setminus T$ and $-1$ on $T\setminus S$.  The active sites form a
uniform $2m$-subset of $[2n]$, independently of a uniform ordering of $m$
plus and $m$ minus signs.  To see this directly, for every fixed signed active
configuration the common intersection $S\cap T$ can be any $(n-m)$-subset of
the remaining $2n-2m$ sites, a number independent of that configuration.

Write the active locations as $J_1<\cdots<J_{2m}$ and let $W_j$ be the sign
partial sum after the $j$th active location.  Lemma~\ref{lem:transport} gives
\begin{equation}\label{eq:stretched-area}
 D_{2n,n}=\sum_{j=1}^{2m-1}|W_j|(J_{j+1}-J_j).
\end{equation}
The $2m+1$ gaps before, between, and after the active sites are exchangeable
and contain $2n-2m$ inactive sites in total.  Indeed, the map
\[
 (J_1,\ldots,J_{2m})\longmapsto
 (J_1-1,J_2-J_1-1,\ldots,J_{2m}-J_{2m-1}-1,2n-J_{2m})
\]
is a bijection from active $2m$-subsets to weak compositions of $2n-2m$
into $2m+1$ parts.  The uniform law on active subsets therefore makes the
gaps exchangeable.  Consequently,
\[
 \E(J_{j+1}-J_j\mid M=m)=1+\frac{2n-2m}{2m+1}
 =\frac{2n+1}{2m+1}.
\]
The active locations and signs are independent, so
Lemma~\ref{lem:discrete-bridge-area} yields
\[
 \E(D_{2n,n}\mid M=m)
 =\frac{2n+1}{2m+1}\frac{m4^m}{2\binom{2m}{m}}.
\]
For $m=0$ the distance is zero.  Averaging the displayed identity for
$1\leq m\leq n$ with \eqref{eq:M-law} therefore proves
\eqref{eq:central-exact}.
\end{proof}

\subsection{Central finite-size expansion}

\begin{theorem}[Central finite-size correction from the closed form]\label{thm:finite-size}
The central-component mean has the expansion
\begin{equation}\label{eq:finite-size}
 \E D_{2n,n}
 =\frac{\sqrt{2\pi}}4n^{3/2}
 -\frac{5\sqrt{2\pi}}{64}n^{1/2}
 +O(n^{-1/2}).
\end{equation}
\end{theorem}

\begin{proof}
The Stirling series and its positive-real remainder bound, applied to
$\log(m!)=\log m+\log\Gamma(m)$,
\cite[\href{https://dlmf.nist.gov/5.11.E1}{Eq.~(5.11.1)} and
\href{https://dlmf.nist.gov/5.11.ii}{Section~5.11(ii)}]{NISTDLMF} give
\[
 \log(m!)=\left(m+\frac12\right)\log m-m+\frac12\log(2\pi)
 +\frac1{12m}-\frac1{360m^3}+O(m^{-5})
\]
and hence, as $m\to\infty$,
\[
 \binom{2m}{m}
 =\frac{4^m}{\sqrt{\pi m}}
  \left(1-\frac1{8m}+\frac1{128m^2}+O(m^{-3})\right).
\]
Write $C_m=\binom{2m}{m}$.  Applying this expansion at $m=n$ and
$m=2n+1$ yields
\begin{align*}
 \frac{n^2}{8n+2}
 &=\frac n8\left(1-\frac1{4n}+O(n^{-2})\right),\\
 \frac{C_{2n+1}}{C_n^2}
 &=\frac{4\pi n}{\sqrt{\pi(2n+1)}}
   \frac{1-1/(8(2n+1))+O(n^{-2})}
        {1-1/(4n)+O(n^{-2})}\\
 &=2\sqrt{2\pi n}
   \left(1-\frac1{4n}+O(n^{-2})\right)
   \left(1-\frac1{16n}+O(n^{-2})\right)
   \left(1+\frac1{4n}+O(n^{-2})\right)\\
 &=2\sqrt{2\pi n}
   \left(1-\frac1{16n}+O(n^{-2})\right).
\end{align*}
Substitution in the published closed form
\eqref{eq:central-closed-form} gives
\[
 \E D_{2n,n}
 =\frac{\sqrt{2\pi}}4n^{3/2}
  \left(1-\frac5{16n}+O(n^{-2})\right),
\]
which is \eqref{eq:finite-size}.
\end{proof}

\subsection{Exact finite-sum formula and numerical evaluation}

The transport identity gives an exact one-dimensional computation that avoids
enumerating pairs of subsets.  Let $H_{N,k,r}$ count how many sites of a
uniform $k$-subset of $[N]$ lie in a fixed marked set of size $r$.  Explicitly,
\[
 \Pp(H_{N,k,r}=j)
 =\frac{\binom rj\binom{N-r}{k-j}}{\binom Nk},
 \qquad
 \max\{0,k-N+r\}\leq j\leq\min\{k,r\}.
\]
Let $H'_{N,k,r}$ be an independent copy.  Then
\begin{proposition}[Exact hypergeometric sum]\label{prop:exact-computation}
\[
 \E D_{N,k}=\sum_{r=1}^{N-1}
 \E|H_{N,k,r}-H'_{N,k,r}|.
\]
For a probability mass function $p_j=\Pp(H_{N,k,r}=j)$, each summand admits
the linear-time evaluation
\[
 \E|H-H'|
 =2\sum_jp_j\left(j\sum_{h<j}p_h-\sum_{h<j}hp_h\right).
\]
\end{proposition}

\begin{proof}
The first identity follows from Lemma~\ref{lem:transport}, Tonelli's theorem,
and the fact that $C_S(r)$ has the stated hypergeometric law.  For the second,
split the double sum $\sum_{h,j}|j-h|p_jp_h$ into the two regions $h<j$ and
$h>j$ and use symmetry.
\end{proof}

The OpenMP program \texttt{code/exact\_mean.cpp} numerically evaluates the
exact finite sum in Proposition~\ref{prop:exact-computation} with normalized
log-binomial weights.
It reports $\E D_{2n,k}/n^{3/2}$ together with the constant in
Theorem~\ref{thm:main} and the finite-size residuals.  In the central case it
also evaluates Proposition~\ref{prop:central-exact} independently.
Tables~\ref{tab:numerics} and~\ref{tab:noncentral} show representative
outputs; the predicted
coefficient of $n^{1/2}$ in Theorem~\ref{thm:finite-size} is
$-5\sqrt{2\pi}/64=-0.195830333955547\ldots$.
Here and in Table~\ref{tab:numerics},
$c=\sqrt{2\pi}/4=0.626657068657750\ldots$.

\begin{table}[hbp]
\centering
\small
\begin{tabular}{rccc}
$n$ & $\E D_{2n,n}/n^{3/2}$
& $(\E D_{2n,n}-cn^{3/2})/\sqrt n$
& \shortstack{absolute formula\\difference}\\
\hline
25   & 0.618994383 & -0.191567130 & $<10^{-13}$\\
100  & 0.624709599 & -0.194746989 & $7\times10^{-14}$\\
400  & 0.626168173 & -0.195558381 & $9\times10^{-13}$\\
800  & 0.626412451 & -0.195694264 & $2\times10^{-11}$\\
1600 & 0.626534717 & -0.195762276 & $4\times10^{-11}$
\end{tabular}
\caption{Central-component floating-point evaluations.  The last column is
the absolute difference between the symmetric-difference sum in
Proposition~\ref{prop:central-exact} and the hypergeometric sum in
Proposition~\ref{prop:exact-computation}.}
\label{tab:numerics}
\end{table}

\begin{table}[hbp]
\centering
\small
\begin{tabular}{cccr@{\qquad}r}
$\rho$ & $n$ & $\kappa$ & $\E D_{2n,\kappa}/n^{3/2}$
& limiting constant\\
\hline
$1/4$  & 100 & 50  & 0.539679868 & 0.542700941\\
$1/4$  & 800 & 400 & 0.542319853 & 0.542700941\\
$1/10$ & 100 & 20  & 0.369966930 & 0.375994241\\
$1/10$ & 800 & 160 & 0.375223284 & 0.375994241
\end{tabular}
\caption{Noncentral evaluations of the hypergeometric sum in
Proposition~\ref{prop:exact-computation}.  The last column is the corresponding
constant $\frac{\sqrt{2\pi}}2\sqrt{\rho(1-\rho)}$ from
Theorem~\ref{thm:main}.}
\label{tab:noncentral}
\end{table}
The values in both tables are direct floating-point evaluations of
deterministic finite sums, not Monte Carlo estimates or certified interval
bounds.

\section{Conclusion}

The matchgate components of West et al. are cover graphs of rectangular
minuscule lattices, so
the distance law and moment convergence of Defant et al. imply the
$n^{3/2}$ central growth in Conjecture~14.  The cumulative-count argument gives
an independent finite-population proof, including arbitrary density sequences
and explicit moment bounds.  The minuscule-lattice Wiener formula of Defant et
al., also recorded in discrete earth-mover notation by Bourn and Erickson,
supplies the closed-form finite-size mean and its leading $n^{1/2}$ correction.
The central conditional-on-$|S\setminus T|$ decomposition gives a different
finite-sum identity.
For complexity, the distance asymptotic implies
Corollary~\ref{cor:complexity} only after averaging both the initial Pauli
vertex and the Haar matchgate unitary.  Corollary~\ref{cor:krylov} is a separate
pointwise comparison for the specified Hamiltonian
\eqref{eq:matchgate-Hamiltonian}; it does not identify a Haar-averaged Krylov
quantity without an additional dynamical law.

\section*{AI-use disclosure}

OpenAI GPT-5.6 Sol (model identifier \texttt{gpt-5.6-sol}) served as an
interactive research assistant for developing and stress-testing arguments,
auditing citations and calculations, preparing exact verification code, and
revising the manuscript.  The author directed the research, selected and
checked the final arguments, and takes responsibility for all mathematical
claims, citations, computations, and wording.

\section*{Data and code availability}

The source archive accompanying this manuscript contains the checker and its
recorded output under \texttt{code/} and \texttt{results/}, together with the
\texttt{Makefile}.  The versioned source package is available at
\url{https://github.com/skypher/matchgate}.  The manuscript version
corresponding to this source package is archived at Git tag
\href{https://github.com/skypher/matchgate/tree/paper-2026-09-04-r6}%
{\texttt{paper-2026-09-04-r6}}.  It is built and run, for example, with
\begin{gather*}
 \texttt{make paper},\qquad \texttt{make exact-mean},\\
 \texttt{./build/exact\_mean 100 400 800 1600},\\
 \texttt{make check-exact-mean}.
\end{gather*}
The manuscript build requires pdf\LaTeX.  The numerical build requires a C++20
compiler with OpenMP support, and the checker requires Python~3.10 or later.
The values in
Tables~\ref{tab:numerics} and~\ref{tab:noncentral} were produced with
GCC~13.3.0.  The implementation
uses \texttt{long double}, \texttt{lgammal}, and normalized exponential
weights to evaluate the exact finite-sum identities numerically; reported
quantities are printed directly at \texttt{long double} round-trip precision.
The \texttt{check-exact-mean} target first checks the token-graph metric
assertion in Proposition~\ref{prop:graph-metric}, Propositions
\ref{prop:closed-form} and \ref{prop:exact-computation}, and
Lemma~\ref{lem:transport} by exact rational enumeration for $N\leq8$,
Lemma~\ref{lem:discrete-bridge-area} for $m\leq6$, and
Proposition~\ref{prop:central-exact} for $n\leq5$.  It then replays every
recorded floating-point row and fails if an output leaves its stated numerical
tolerance.

\begin{thebibliography}{99}
\footnotesize\raggedright

\bibitem{AudenaertEtAl}
K. M. R. Audenaert, C. Godsil, G. Royle, and T. Rudolph,
\emph{Symmetric squares of graphs},
Journal of Combinatorial Theory, Series B \textbf{97} (2007), no.~1, 74--90.
\href{https://doi.org/10.1016/j.jctb.2006.04.002}%
{doi:10.1016/j.jctb.2006.04.002}.

\bibitem{BournErickson}
R. Bourn and W. Q. Erickson,
\emph{Palindromicity of the numerator of a statistical generating function},
Discrete Mathematics \textbf{348} (2025), no.~3, 114336.
\href{https://doi.org/10.1016/j.disc.2024.114336}%
{doi:10.1016/j.disc.2024.114336}.

\bibitem{BournWillenbring}
R. Bourn and J. F. Willenbring,
\emph{Expected value of the one-dimensional earth mover's distance},
Algebraic Statistics \textbf{11} (2020), no.~1, 53--78.
\href{https://doi.org/10.2140/astat.2020.11.53}%
{doi:10.2140/astat.2020.11.53}.

\bibitem{DeAngelisGray}
M. De Angelis and A. Gray,
\emph{Why the 1-Wasserstein distance is the area between the two marginal
CDFs}, arXiv:2111.03570 (2021).
\href{https://doi.org/10.48550/arXiv.2111.03570}{doi:10.48550/arXiv.2111.03570}.

\bibitem{DefantEtAl}
C. Defant, V. F\'eray, P. Nadeau, and N. Williams,
\emph{Wiener indices of minuscule lattices},
The Electronic Journal of Combinatorics \textbf{31} (2024), no.~1,
Paper No.~P1.41.
\href{https://doi.org/10.37236/12002}{doi:10.37236/12002}.

\bibitem{delBarrioGineMatran}
E. del Barrio, E. Gin\'e, and C. Matr\'an,
\emph{Central limit theorems for the Wasserstein distance between the empirical
and the true distributions},
Annals of Probability \textbf{27} (1999), no.~2, 1009--1071.
\href{https://doi.org/10.1214/aop/1022677394}{doi:10.1214/aop/1022677394}.

\bibitem{delBarrioCorrection}
E. del Barrio, E. Gin\'e, and C. Matr\'an,
\emph{Correction: Central limit theorems for the Wasserstein distance between
the empirical and the true distributions},
Annals of Probability \textbf{31} (2003), no.~2, 1142--1143.
\href{https://projecteuclid.org/journals/annals-of-probability/volume-31/issue-2/Correction--Central-limit-theorems-for-the-Wasserstein-distance-between/aop/1048516548.pdf}%
{Project Euclid: aop/1048516548}.

\bibitem{DiazEtAl}
N. L. Diaz, D. Garc\'ia-Mart\'in, S. Kazi, M. Larocca, and M. Cerezo,
\emph{Showcasing a barren plateau theory beyond the dynamical Lie algebra},
arXiv:2310.11505 (2023).
\href{https://doi.org/10.48550/arXiv.2310.11505}%
{doi:10.48550/arXiv.2310.11505}.

\bibitem{FabilaMonroy}
R. Fabila-Monroy, D. Flores-Pe\~naloza, C. Huemer, F. Hurtado, J. Urrutia,
and D. R. Wood,
\emph{Token graphs},
Graphs and Combinatorics \textbf{28} (2012), no.~3, 365--380.
\href{https://doi.org/10.1007/s00373-011-1055-9}%
{doi:10.1007/s00373-011-1055-9}.

\bibitem{Hoeffding}
W. Hoeffding,
\emph{Probability inequalities for sums of bounded random variables},
Journal of the American Statistical Association \textbf{58} (1963), no.~301,
13--30.
\href{https://doi.org/10.1080/01621459.1963.10500830}%
{doi:10.1080/01621459.1963.10500830}.

\bibitem{IbarraRivera}
S. Ibarra and L. M. Rivera,
\emph{The automorphism groups of some token graphs},
Proyecciones Journal of Mathematics \textbf{42} (2023), no.~6, 1627--1651.
\href{https://doi.org/10.22199/issn.0717-6279-5954}%
{doi:10.22199/issn.0717-6279-5954}.

\bibitem{Janson}
S. Janson,
\emph{Brownian excursion area, Wright's constants in graph enumeration, and
other Brownian areas},
Probability Surveys \textbf{4} (2007), 80--145.
\href{https://doi.org/10.1214/07-PS104}{doi:10.1214/07-PS104}.

\bibitem{NISTDLMF}
\emph{NIST Digital Library of Mathematical Functions},
\url{https://dlmf.nist.gov/}, Release 1.2.7 of June 15, 2026.
F. W. J. Olver, A. B. Olde Daalhuis, D. W. Lozier, B. I. Schneider,
R. F. Boisvert, C. W. Clark, B. R. Miller, B. V. Saunders, H. S. Cohl,
and M. A. McClain, eds.

\bibitem{ParkerEtAl}
D. E. Parker, X. Cao, A. Avdoshkin, T. Scaffidi, and E. Altman,
\emph{A universal operator growth hypothesis},
Physical Review X \textbf{9} (2019), no.~4, 041017.
\href{https://doi.org/10.1103/PhysRevX.9.041017}%
{doi:10.1103/PhysRevX.9.041017}.

\bibitem{Rosen}
B. Ros\'en,
\emph{Limit theorems for sampling from finite populations},
Arkiv f\"or Matematik \textbf{5} (1965), no.~5, 383--424.
\href{https://doi.org/10.1007/BF02591138}{doi:10.1007/BF02591138}.

\bibitem{Schwerdtfeger}
U. Schwerdtfeger,
\emph{Linear functional equations with a catalytic variable and area limit
laws for lattice paths and polygons},
European Journal of Combinatorics \textbf{36} (2014), 608--640.
\href{https://doi.org/10.1016/j.ejc.2013.10.004}{doi:10.1016/j.ejc.2013.10.004}.

\bibitem{Vallender}
S. S. Vallender,
\emph{Calculation of the Wasserstein distance between probability
distributions on the line},
Theory of Probability and Its Applications \textbf{18} (1974), no.~4,
784--786.
\href{https://doi.org/10.1137/1118101}{doi:10.1137/1118101}.

\bibitem{VallenderAddendum}
S. S. Vallender,
\emph{Addendum: Calculation of the Wasserstein distance between probability
distributions on the line},
Theory of Probability and Its Applications \textbf{26} (1982), no.~2, 435.
\href{https://doi.org/10.1137/1126051}{doi:10.1137/1126051}.

\bibitem{WestEtAl}
M. West, N. Dowling, A. Southwell, M. Sevior, M. Usman, K. Modi, and T. Quella,
\emph{A graph-theoretic approach to chaos and complexity in quantum systems},
SciPost Physics Core \textbf{8} (2025), no.~4, 081.
\href{https://doi.org/10.21468/SciPostPhysCore.8.4.081}%
{doi:10.21468/SciPostPhysCore.8.4.081}.

\end{thebibliography}

\end{document}
